\clearpage\pagestyle{otMain}\fontsize{11}{15}\selectfont
\corrigesChapitre{3}{Calcul numérique et littéral}
\begin{blocLivre}{cor-3-1}
\begin{corrige}{1}{Puissances et parenthèses}
\demarche{Lire la base avant d’appliquer une règle de puissance.}
Les parenthèses font porter la puissance sur le signe : $(-2)^4=\res{16}$ tandis que $-2^4=-(2^4)=\res{-16}$. Un exposant négatif signifie un inverse : $3^{-2}=\res{1/9}$. Enfin $5^0=\res1$, car toute base non nulle à la puissance 0 vaut 1.
\end{corrige}
\end{blocLivre}
\begin{blocLivre}{cor-3-2}
\begin{corrige}{2}{Racines simplifiées}
\demarche{Décomposer le radicande en carré parfait multiplié par un reste.}
On extrait le plus grand carré parfait : $\sqrt{50}=\sqrt{25\times2}=\res{5\sqrt2}$ ; $\sqrt{108}=\sqrt{36\times3}=\res{6\sqrt3}$. Pour le quotient, $\sqrt{16/25}=4/5$, les racines du numérateur et du dénominateur étant positives.
\end{corrige}
\end{blocLivre}
\begin{blocLivre}{cor-3-3}
\begin{corrige}{3}{Écriture scientifique}
\demarche{Regrouper les coefficients puis additionner ou soustraire les exposants.}
On sépare les coefficients et les puissances : $A=(3\times4/6)\times10^{5-2-2}=2\times10^1$. Cette écriture est scientifique puisque $1\le2<10$. On obtient $\res{A=2\times10^1=20}$.
\end{corrige}
\end{blocLivre}
\begin{blocLivre}{cor-3-4}
\begin{corrige}{4}{Développer et réduire}
\demarche{Effectuer toutes les distributions avant de regrouper les termes semblables.}
Pour $A$, distribuer le signe négatif à tous les termes : $A=3x-6-4x-10=\res{-x-16}$. Pour $B$, chaque terme du premier facteur multiplie chaque terme du second : $B=2x^2-x+8x-4=\res{2x^2+7x-4}$.
\end{corrige}
\end{blocLivre}
\begin{blocLivre}{cor-3-5}
\begin{corrige}{5}{Reconnaître une identité}
\demarche{Identifier les deux termes de l’identité remarquable.}
La première expression vaut $\res{4x^2-12x+9}$. La différence de carrés se factorise : $x^2-25=\res{(x-5)(x+5)}$. Enfin $x^2+6x+9=x^2+2\times x\times3+3^2=\res{(x+3)^2}$. Les identités s’utilisent dans les deux sens.
\end{corrige}
\end{blocLivre}
\begin{blocLivre}{cor-3-6}
\begin{corrige}{6}{Un facteur commun}
\demarche{Repérer un facteur entier commun, même s’il contient plusieurs termes.}
Dans $A$, le facteur commun est $x-1$ : $A=\res{(x-1)(5x+3)}$. Dans $B$, extraire $x+2$ donne $(x+2)((x+2)-4)$, donc $B=\res{(x+2)(x-2)}$. Développer les résultats permet de vérifier que rien n’a été perdu.
\end{corrige}
\end{blocLivre}
\begin{blocLivre}{cor-3-7}
\begin{corrige}{7}{Une simplification conditionnelle}
\demarche{Déterminer le domaine avant de simplifier.}
Le dénominateur impose $x\ne-2$. Pour ces valeurs, factoriser $x^2-4=(x-2)(x+2)$ puis simplifier le facteur non nul : l’expression vaut $\res{x-2}$. Pour $x=3$, elle vaut 1, ce que confirme $(9-4)/(3+2)=1$. La valeur $-2$ reste interdite.
\end{corrige}
\end{blocLivre}
\begin{blocLivre}{cor-3-8}
\begin{corrige}{8}{Addition de quotients}
\demarche{Construire un dénominateur commun, puis additionner les numérateurs.}
Le dénominateur commun est $3x$. Ainsi $2/x=6/(3x)$ et la somme vaut $(6+x-1)/(3x)=\res{(x+5)/(3x)}$. Le domaine reste $x\ne0$. On n’additionne pas les dénominateurs : ils décrivent la même taille d’unité après mise au même dénominateur.
\end{corrige}
\end{blocLivre}
\begin{blocLivre}{cor-3-9}
\begin{corrige}{9}{Isoler une grandeur}
\demarche{Utiliser les opérations inverses et le domaine de la grandeur.}
La première relation donne $\res{t=d/v}$. Pour la deuxième, $r^2=S/\pi$, puis la condition $r\ge0$ sélectionne $\res{r=\sqrt{S/\pi}}$. Enfin $-2y=7-3x$, donc $\res{y=(3x-7)/2}$. Le signe moins devant $2y$ doit être conservé jusqu’à la division.
\end{corrige}
\end{blocLivre}
\begin{blocLivre}{cor-3-10}
\begin{corrige}{10}{Comparer par différence ou par quotient}
\demarche{Ramener une comparaison au signe d’une différence ou à la position d’un quotient par rapport à 1.}
On calcule $A-B=x^2+2x+1-x^2-1=2x$. Comme $x\ge0$, cette différence est positive ou nulle. Ainsi $\res{A\ge B}$, avec égalité exactement lorsque $2x=0$, soit $\res{x=0}$. Ce raisonnement vaut pour toutes les valeurs autorisées de $x$. Les deux fractions sont strictement positives et $\dfrac79\div\dfrac9{11}=\dfrac{77}{81}<1$ : donc $\res{\dfrac79<\dfrac9{11}}$. Le premier nombre vaut environ 0,95 fois le second.
\end{corrige}
\end{blocLivre}
\begin{blocLivre}{cor-3-11}
\begin{corrige}{11}{Un programme de calcul}
\demarche{Transformer le programme en expression, puis exploiter la parité.}
Le résultat est $(x+3)^2-9=x^2+6x=x(x+6)$. Si $x=2k$, il devient $2k(2k+6)=4k(k+3)$. Comme $k(k+3)$ est entier, le résultat est multiple de 4. La factorisation $\res{x(x+6)}$ rend cette propriété plus visible que la description du programme.
\end{corrige}
\end{blocLivre}
\begin{blocLivre}{cor-3-12}
\begin{corrige}{12}{Calculer astucieusement}
\demarche{Choisir l’identité qui correspond directement à la forme donnée.}
La différence de carrés donne $(1001-999)(1001+999)=2\times2000=\res{4000}$. En général, $(n+1)^2-(n-1)^2=2\times2n=\res{4n}$. On aurait le même résultat en développant, mais la factorisation réduit ici le nombre d’opérations.
\end{corrige}
\end{blocLivre}
\begin{blocLivre}{cor-3-13}
\begin{corrige}{13}{Un carré encadré}
\demarche{Soustraire les aires et compter les deux largeurs de bordure.}
Le grand carré a pour côté $x+4$, car la bordure s’ajoute des deux côtés. L’aire cherchée est $(x+4)^2-x^2=8x+16=8(x+2)$. Résoudre $8x+16=56$ donne $\res{x=5}$. On contrôle : $9^2-5^2=81-25=56$.
\label{fig-cor-3-13}
\figureGeo[corrige]{bordure}
\end{corrige}
\end{blocLivre}
\begin{blocLivre}{cor-3-14}
\begin{corrige}{14}{Une identité trompeuse}
\demarche{Résoudre l’égalité obtenue sans accepter la transformation erronée.}
Il faut d’abord $x\ne-2$. L’égalité à 2 équivaut alors à $x+4=2(x+2)$, soit $x+4=2x+4$, donc $\res{x=0}$. Cette valeur est autorisée et donne $4/2=2$. Une égalité vraie pour une seule valeur ne justifie pas la fausse règle de simplification utilisée.
\end{corrige}
\end{blocLivre}
\begin{blocLivre}{cor-3-15}
\begin{corrige}{15}{Deux modèles équivalents}
\demarche{Transformer l’expression avant de décider du degré de l’équation.}
La différence de carrés se factorise : $A(x)=((x+5)-(x-1))((x+5)+(x-1))=6(2x+4)=12x+24$. Les termes en $x^2$ s’annulent. L’équation devient $12x+24=72$, d’où $\res{x=4}$. La condition $x\ge1$ est respectée et $9^2-3^2=72$.
\end{corrige}
\end{blocLivre}
